% ============================================================
% Section 163: Morse 势的 s 态能级
% ============================================================
\section{量子力学：Morse 势的 \texorpdfstring{$s$}{s} 态能级}
\label{sec:morse-potential-s-level}

\begin{modelbox}[题目：Morse 势的 \texorpdfstring{$s$}{s} 态能级]
双原子分子中两原子间的作用可用 Morse 势
\[
V(r)=V_0\bigl[1-e^{-(r-R)/a}\bigr]^2-V_0
\]
描写。对分子内部相对运动的 $s$ 态（$l=0$），解出径向方程并确定束缚态能级公式。
\end{modelbox}

\begin{tipbox}[考点快速分析]
\begin{itemize}
    \item \textbf{Morse 势}：$V(r\to\infty)\to0$，$V(R)=-V_0$（极小值），$a$ 为势阱宽度参数
    \item \textbf{变量代换}：$\xi=e^{-(r-R)/a}$，将微分方程化为合流超几何方程
    \item \textbf{束缚态条件}：合流超几何函数退化为多项式（截断条件）
    \item \textbf{物理意义}：谐振子项 $+$ 非谐修正项
\end{itemize}
\end{tipbox}

\begin{derivationbox}[标准解题过程]
\textbf{1. 径向方程与变量代换}

$s$ 态波函数 $\psi(r)=u(r)/r$，$u(r)$ 满足
\[
-\frac{\hbar^2}{2\mu}\frac{d^2u}{dr^2}+V(r)u=Eu,\quad -V_0<E<0.
\]

令 $\xi=e^{-(r-R)/a}$，则 $\dfrac{d}{dr}=-\dfrac{\xi}{a}\dfrac{d}{d\xi}$。代入整理得
\begin{equation}
\xi^2\frac{d^2u}{d\xi^2}+\xi\frac{du}{d\xi}+\bigl(2\alpha^2\xi-\alpha^2\xi^2-\beta^2\bigr)u=0,
\end{equation}
其中
\[
\alpha=\frac{a}{\hbar}\sqrt{2\mu V_0},\quad \beta=\frac{a}{\hbar}\sqrt{-2\mu E}.
\]

\textbf{2. 渐近行为与标准形式}

$\xi\to0$（$r\to\infty$）：$u\sim\xi^\beta$（取衰减解）；
$\xi\to\infty$（$r\to0$）：$u\sim e^{-\alpha\xi}$（取衰减解）。

设 $u(\xi)=\xi^\beta e^{-\alpha\xi}F(\xi)$，代入得 $F$ 满足
\[
\xi F''+(2\beta+1-2\alpha\xi)F'-\bigl(\alpha+2\alpha\beta-2\alpha^2\bigr)F=0.
\]

再令 $\eta=2\alpha\xi$，化为标准 Kummer 方程：
\[
\eta F''+(2\beta+1-\eta)F'-\Bigl(\beta+\frac12-\alpha\Bigr)F=0.
\]

\textbf{3. 束缚态条件与能级}

正则解为 $_1F_1(\beta+\frac12-\alpha,\,2\beta+1;\,\eta)$。为保证 $\eta\to\infty$ 时 $u\to0$，合流超几何函数必须退化为多项式，要求
\[
\beta+\frac12-\alpha=-\nu,\quad \nu=0,1,2,\ldots
\]
即 $\beta=\alpha-(\nu+\frac12)$。代入 $\beta$ 的定义：
\begin{equation}
\boxed{E_\nu=-\frac{\hbar^2}{2\mu a^2}\Bigl[\alpha-\Bigl(\nu+\frac12\Bigr)\Bigr]^2,\quad \nu=0,1,2,\ldots}
\end{equation}

展开并引入特征频率 $\omega=\sqrt{\dfrac{2V_0}{\mu a^2}}$：
\begin{equation}
\boxed{E_\nu=-V_0+\Bigl(\nu+\frac12\Bigr)\hbar\omega-\Bigl(\nu+\frac12\Bigr)^2\frac{\hbar^2}{2\mu a^2}.}
\end{equation}

\textbf{4. 适用条件}

要求 $\beta>0$，即 $\alpha>\nu+\frac12$，故能级数有限。至少存在一条束缚态要求
\[
V_0a^2>\frac{\hbar^2}{32\mu}.
\]
\end{derivationbox}

\begin{formulabox}[答案汇总]
\begin{center}
\begin{tabular}{ll}
\toprule
\textbf{结论} & \textbf{结果} \\
\midrule
束缚态能级 & $E_\nu=-V_0+\left(\nu+\frac12\right)\hbar\omega-\left(\nu+\frac12\right)^2\dfrac{\hbar^2}{2\mu a^2}$ \\
参数 & $\omega=\sqrt{2V_0/(\mu a^2)}$，$\nu=0,1,2,\ldots$ \\
束缚态数 & 有限，满足 $\nu+\frac12<\alpha=\dfrac{a}{\hbar}\sqrt{2\mu V_0}$ \\
\bottomrule
\end{tabular}
\end{center}
\end{formulabox}

\begin{insightbox}[物理图像]
\begin{itemize}
    \item Morse 势是能精确求解的双原子分子势模型。能级公式中第一项 $-V_0$ 是解离能；第二项为谐振子近似；第三项为负的非谐修正，使能级间距随 $\nu$ 增大而减小。
    \item 这一能级结构特征与真实分子振动光谱（如 HCl、H$_2$）高度吻合：基频附近近似等间距，高激发时向解离限收敛。
    \item 与谐振子不同，Morse 势的束缚态数量有限，反映了真实分子只能存在有限个振动激发态的事实。
\end{itemize}
\end{insightbox}
